
%% Definition Uni-Logo und Linie darunter
%
\newcommand{\UniLogo}{%
\raisebox{-1.0cm}{%
\fboxsep0pt%
\hyperlink{IHV}{%
\includegraphics[viewport= 0 0 145 175,height=2.7cm]{L-ALU-ab-2009-farbig}}}
\\[-1.25cm]
\rule{13.5cm}{2pt}}%
%
\newcommand{\UniLogoFront}{%
\raisebox{-1.8cm}{%
\fboxsep0pt{%
%\includegraphics[viewport= 0 0 2430 2880,height=4.7cm]{L-ALU}}}
\includegraphics[viewport= 0 0 243 288,height=4.7cm]{L-ALU-ab-2009-farbig}}}
\\[-1.6cm]
\rule{12.0cm}{3pt}}
%
\newcommand{\UniLogoMathFak}{%
\raisebox{-1.0cm}{%
\fboxsep0pt%
\hyperlink{IHV}{%
\includegraphics[viewport= 0 0 145 175,height=2.7cm]{L-ALU-ab-2009-farbig}}}
\\[-.95cm]
\rule{13.5cm}{2pt}}%


%% Definition der Koepfe %%
%
\newcommand{\newllogo}[5]{\expandafter\newcommand\csname#1\endcsname{{%
#2%
\begin{minipage}[b]{#3}%
\large\textbf{\textsf{#4}\mbox{}}
\end{minipage} \hfill %
\Semester\qquad#5\vspace{10mm}}}
}
%
\newcommand{\newlogo}[4]{\newllogo{#1}{#2}{#3}{#4}{\UniLogo}}
\newcommand{\newmlogo}[4]{\newllogo{#1}{#2}{#3}{#4}{\UniLogoMathFak}}

% Abteilung f"ur Angewandte Mathematik (\iam)
\newlogo{iam}{%
\includegraphics[height=1.5cm]{L-AM}\vspace{2ex}
\hspace*{1.3cm} % <- Dieses Leerzeichen ist notwendig...
}{5cm}{Abteilung f{\"u}r\\Angewandte Mathematik\\[-0.7em]}
%
% Abteilung f"ur Mathematische Logik (\ifl)
\newlogo{ifl}{%
\includegraphics[height=1.5cm]{L-Logik}\vspace{2ex}
\hspace*{2.3cm} % <- hier auch...
}{4.5cm}{Abteilung f{\"u}r\\ Mathematische Logik\\[-.7em]}
%
% Abteilung f"ur Mathematische Stochastik (\ifs)
\newlogo{ifs}{%
\includegraphics[height=1.5cm]{L-MS}\vspace{2ex}
\hfill
}{3cm}{}
%
% Abteilung f{\"u}r Reine Mathematik (\mai)
\newlogo{mai}{%
\rule{0pt}{1.5cm}\vspace{2ex}%
}{5cm}{%
Abteilung f{\"u}r\\ Reine Mathematik\\[-.7em]}
%
% Abteilung f{\"u}r Didaktik der Mathematik (\mad)
\newlogo{mad}{%
\rule{0pt}{1.5cm}\vspace{2ex}%
}{5cm}{%
Abteilung f{\"u}r\\Didaktik der Mathematik\\[-0.7em]}
%
% Didaktik -- PH Freiburg (madph)
\newlogo{madph}{%
\includegraphics[height=1.6cm]{Definition/L-PH}\hspace{3.0cm}\includegraphics[height=1.5cm]{Definition/L-PH-FACE}\vspace{2ex}%
}{1cm}{%
\\[-0.7em]}
%
% Abteilung f{\"u}r Finanzmathematik (\fin)
\newlogo{fin}{%
\rule{0pt}{1.5cm}\vspace{2ex}%
}{5cm}{%
Abteilung f{\"u}r\\Finanzmathematik\\[-0.7em]}
%
% Abteilung f{\"u}r Quantitative Finanzmarktforschung (\qff)
\newlogo{qff}{%
\rule{0pt}{1.5cm}%
\vspace{2ex}%
% \hspace*{4.0cm}%
}{7.5cm}{%
Abteilung f{\"u}r\\Quantitative Finanzmarktforschung\\[-0.7em]}
%
% Abteilung Medizinische Biometrie und Statistik am
% Institut f"ur Med. Biometrie und Med. Informatik (\imb)
\newlogo{imb}{%
\includegraphics[height=1.5cm]{L-imbi}\vspace{2ex}%
\hfill%
}{5.6cm}{%
Institut f{\"u}r\\Medizinische Biometrie und\\
Statistik\\[-0.7em]}
%
% Computer Vision Group (\cvg)
\newlogo{cvg}{%
\raisebox{0.8ex}{\includegraphics[height=.5cm]{Definition/L-Vision-2}}\vspace{2ex}
\hspace*{1.3cm} % <- Dieses Leerzeichen ist notwendig...
}{5cm}{Computer Vision Group\\[-0.7em]}
%
% Vorsitzender Pruefungsausschuesse (\pruef)
\newmlogo{pruef}{}{10cm}{%
Mathematisches Institut\\
Vorsitzender der 
Pr\"ufungsaussch\"usse Mathematik\\
Prof. Dr. S. Goette\\[-.7em]}
%
%

% Math. Inst. (\minst)
\newlogo{minst}{%
\rule{0pt}{1.5cm}\vspace{2ex}%
}{5cm}{%
Mathematisches Institut\\[-.7em]}
%

% Mathematisches Institut (\fak)
\newmlogo{fak}{}{5cm}{%
Mathematisches Institut\\[-.7em]}
%
% Mathematische und Physikalische Fakult"at (\mfak)
\newlogo{mfak}{%
\rule{0pt}{2.5cm}%
}{6.3cm}{%
\begin{center}%
Fakult\"at f\"ur\\%
Mathematik und Physik\\%
Mathematisches Institut\\%
\end{center}}
%
% Leere Kopfzeile mit Unilogo   (\leer)
\newmlogo{leer}{%
\rule{0pt}{1.5cm}%
}{10cm}{%
\phantom{Leer g}\\[-.5em]}


%% Definition der Sprechstundenseite %%
%
\newcommand{\Sprechstunden}{%
\fak \\ %
\bigskip\bigskip
{\renewcommand{\baselinestretch}{2.0}\tiny\large

\textbf{\Large Sprechstundenverzeichnis}\\[4ex]
Bitte informieren Sie sich \"uber die aktuell g\"ultigen Sprechstunden\\[2ex] der Professorinnen, der Professoren und Assistent(inn)en auf der Internet-Seite\\[4ex]
\url{http://home.mathematik.uni-freiburg.de/personen/list/sprechstunden.de.html}.}
}


%% Definition des Impressums %%
%
\newcommand{\impressum}{%
\thispagestyle{empty}%
\fakesection{Impressum}%
~\vfill
\hyperlink{IHV}{\textbf{Impressum}}\\[2ex]
Herausgeber: \\[1ex]
Mathematisches Institut \\
Eckerstr. 1\\
79104 Freiburg\\
Tel.: 0761-203-5534\\
E-Mail: institut@math.uni-freiburg.de
}
%
